\section{Results and discussion}
\label{sec:results}

Section~\ref{sec:eda} first characterizes the compiled dataset. The
modeling experiments then answer four research questions. (RQ1)~How
accurately do the models predict RR\% for studied systems at new
operating conditions? (RQ2)~Does this accuracy transfer to systems from
unseen studies? (RQ3)~Can the uncertainty of individual predictions be
quantified reliably under both tasks? (RQ4)~Which factors govern the
predictions, and are they chemically plausible?

\subsection{Dataset characteristics}
\label{sec:eda}

Figure~\ref{fig:eda} summarizes the compiled data. Removal ratios span
the full 0--100\% range with a median above 50\%, but the distribution
differs strongly between pollutants and between LDH metal pairs
(Fig.~\ref{fig:eda}a,b), reflecting both genuine chemistry and the
uneven attention different systems have received in the literature. The
study-size distribution is heavily skewed: the largest single study
contributes 220 of the 1,342 records, the median study only 17, and a
random record-level sample is therefore dominated by a handful of
intensively measured systems. Spearman correlations among the numeric
descriptors (Fig.~S1) show the expected blocks: the Abraham solvation
parameters correlate strongly with molecular weight and with each other,
and the composition-averaged elemental descriptors correlate with the
metal identities they are derived from. This redundancy motivates the
correlated-descriptor experiment in Section~S3 and, more fundamentally,
explains why record-level splits leak: within one study, most descriptor
columns are constant, so any two records of the same study are
near-duplicates in feature space. Missingness is confined to the
synthesis and texture blocks (Table~\ref{tab:stats}): calcination
conditions are unreported for two thirds of the records, and
pH$_{\mathrm{pzc}}$ and pore volume for roughly a quarter, which is why
imputation quality matters and why texture-based mechanistic claims from
compiled data warrant caution.

\begin{figure}[t]
\centering
\includegraphics[width=\linewidth]{fig_eda_rr}
\caption{Removal-ratio distributions by pollutant and by metal pair}
\label{fig:eda}
\end{figure}

\subsection{Within-system interpolation (RQ1)}
\label{sec:rq1}

Table~\ref{tab:interp} answers (RQ1) by comparing the seven models and
the naive baselines under task (T1), reported as mean $\pm$ SD over 20
seeded splits. The tree ensembles lead: XGBoost and HistGB reach test
$R^2$ of $0.84$--$0.85$ with RMSE near 11 percentage points, followed by
Cubist and RF, while LR trails far behind. This is because RR\%
responds nonlinearly to the operating conditions, with thresholds in pH
and dose that additive linear terms cannot represent. The baselines put
these numbers in perspective. The per-study mean, which uses no
descriptors at all, already attains $R^2$ of $0.43$ (MAE 16.7): nearly
half of the variance a record-level split appears to explain is
available from system identity alone, and LR barely improves on it. The
ensembles' further gain to $R^2$ of $0.85$ therefore reflects genuine
modeling of the operating conditions within known systems. The
in-sample $R^2$ of the ensembles approaches 0.99, and the gap to the
test score is the expected overfitting margin of high-capacity learners
rather than a protocol artifact; the cross-validation and test scores
agree closely. Figure~\ref{fig:parity} shows the parity and residual
plots of the best model on the reference split. It is worth noting that
these scores are already 2--3 RMSE points more conservative than a
protocol without duplicate removal on the same data, since deduplication
alone eliminates records whose test values could be looked up verbatim
from training.

\begin{table}[t]
\centering
\caption{Model and baseline performance under within-system
interpolation (T1), mean $\pm$ SD over 20 seeded 80/20 splits
[VALUES FROM repeated\_interp\_summary.csv]}
\label{tab:interp}
\begin{tabular}{lccc}
\toprule
Model & Test $R^2$ & Test MAE & Test RMSE \\
\midrule
XGBoost & $0.84 \pm 0.03$ & $6.3 \pm 0.6$ & $11.1 \pm 1.0$ \\
HistGB  & $0.85 \pm 0.03$ & $6.3 \pm 0.5$ & $10.8 \pm 1.0$ \\
KNN     & $0.73 \pm 0.06$ & $7.8 \pm 0.8$ & $14.4 \pm 1.5$ \\
Cubist  & $0.79 \pm 0.03$ & $8.0 \pm 0.6$ & $12.7 \pm 1.0$ \\
RF      & $0.78 \pm 0.03$ & $8.6 \pm 0.5$ & $13.0 \pm 0.9$ \\
ANN     & $0.54 \pm 0.05$ & $14.4 \pm 0.7$ & $19.1 \pm 0.9$ \\
LR      & $0.40 \pm 0.04$ & $17.4 \pm 0.5$ & $21.7 \pm 0.7$ \\
\midrule
Per-study mean     & $0.43 \pm 0.05$ & $16.7 \pm 0.8$ & $21.2 \pm 0.9$ \\
Per-pollutant mean & $0.18 \pm 0.03$ & $21.4 \pm 0.6$ & $25.4 \pm 0.6$ \\
Global mean        & $-0.01 \pm 0.01$ & $24.9 \pm 0.7$ & $28.1 \pm 0.6$ \\
\bottomrule
\end{tabular}
\end{table}

\subsection{Cross-study generalization (RQ2)}
\label{sec:rq2}

Figure~\ref{fig:loso} answers (RQ2) with the leave-one-study-out audit.
The picture changes sharply: the median per-study MAE of the best models
rises to 14--16 percentage points, two to three times their
interpolation error, and the per-study $R^2$ is positive for only 4 of
the 32 studies large enough to evaluate it. This is because the models
exploit descriptor combinations that identify a material--pollutant
system rather than transferable structure--activity relations; when all
records of a system are withheld, that shortcut disappears. The
baseline comparison adds an important nuance. Descriptor-free lookups
transfer even worse: the global training-mean reaches a median per-study
MAE of 24 points and the per-pollutant mean of 25, so the models do
reduce the naive error by roughly 40 percent, and Cubist beats the
per-pollutant lookup in 41 of the 59 studies. The models therefore
track the level differences between systems while failing to resolve the
condition response within an unseen system, which is exactly why the
per-study $R^2$ (measured against each study's own variance) stays
negative even as the MAE gain over the baselines is large. Screening
applications get a coarse ranking signal from these models, not
engineering precision. The
per-study distribution in Fig.~\ref{fig:loso}b is wide rather than
uniformly poor: for Cubist, 20 of the 59 studies are predicted with
errors comparable to the interpolation regime (MAE below 6 points),
while 9 studies exceed 30 points. Notably, the model ranking also
reorders: Cubist, mid-table under interpolation, attains the lowest
cross-study median, which could be explained by its rule-plus-linear
structure extrapolating more gracefully than deep tree ensembles.

\textbf{Which systems transfer.} Figure~\ref{fig:transfer} breaks the
per-study errors down by chemistry, and the pattern follows literature
density. Ciprofloxacin studies, which dominate the compilation (32 of
59 studies), are predicted with a median MAE of 5 points, and Zn--Al
systems (28 studies) with 5 points, because each held-out study of these
families leaves many close relatives in training. Sparsely covered
chemistries fare far worse: tetracycline, diclofenac, and levofloxacin
studies show median errors above 20 points, and Ni--Fe systems above 30.
Error also grows with the held-out study's own scope: per-study MAE
correlates with both the number of records and the within-study spread
of RR\% ($r \approx 0.45$ for each), meaning studies that explore wide
condition ranges are precisely the ones an outside model cannot resolve.
The learning curve in Fig.~\ref{fig:lc}b is flat between 5 and 40
training studies, suggesting that cross-study error is limited by the
diversity of compiled systems rather than by record count. This
supports our central argument that random-split scores, including those
in Table~\ref{tab:interp}, must not be read as screening skill: the same
models, the same data, and the same preprocessing yield near-zero skill
when the test set contains genuinely new systems.

\begin{figure}[t]
\centering
\includegraphics[width=\linewidth]{fig_parity}
\caption{Parity and residual plots of the best model under interpolation}
\label{fig:parity}
\end{figure}

\begin{figure}[t]
\centering
\includegraphics[width=\linewidth]{fig_loso_audit}
\caption{Leave-one-study-out audit across 59 studies}
\label{fig:loso}
\end{figure}

\begin{figure}[t]
\centering
\includegraphics[width=\linewidth]{fig_learning_curve}
\caption{Interpolation versus unseen-study error and learning curve}
\label{fig:lc}
\end{figure}

\begin{figure}[t]
\centering
\includegraphics[width=\linewidth]{fig_transfer_breakdown}
\caption{Cross-study error by pollutant, metal pair, and study size}
\label{fig:transfer}
\end{figure}

\subsection{Uncertainty and applicability domain (RQ3)}
\label{sec:rq3}

Figure~\ref{fig:conformal} answers (RQ3). Under interpolation, the
split-conformal 90\% intervals of Eq.~(\ref{eq:conformal}) achieve
empirical coverage of 93\% with a half-width of 19 percentage points, so
predictions for studied systems can be certified at a practically useful
precision. Under the new-study task, coverage is maintained (99\%) only
because the study-grouped calibration inflates the half-width to 68
points, an interval spanning most of the 0--100 RR\% range and therefore
carrying little decision value. This is because residual magnitudes
differ systematically between systems, so a single global quantile must
cover the worst of them. The distance-based applicability-domain flag
does not compensate: the mean distance to the five nearest training
records carries essentially no information about realized error in
either task ($r \le 0.06$, Fig.~\ref{fig:conformal}b), although such
distance heuristics are widely used as trust flags. A key finding is
therefore that the uncertainty tools inherit the transfer problem
itself: interpolation predictions can be certified, whereas no
construction we tested produces informative guarantees for unseen
systems, and coverage claims must be validated empirically for the
deployment task.

\begin{figure}[t]
\centering
\includegraphics[width=\linewidth]{fig_conformal_ad}
\caption{Conformal coverage and applicability-domain flag}
\label{fig:conformal}
\end{figure}

\subsection{Feature effects and mechanistic reading (RQ4)}
\label{sec:rq4}

Figure~\ref{fig:shap} and Table~\ref{tab:ranks} answer (RQ4). Adding
the chemistry-informed features leaves both tasks essentially unchanged
(interpolation MAE 5.6 versus 5.8 points; cross-study median 16.4 versus
16.0), indicating that the models already reconstruct these combinations
internally from pH, pH$_{\mathrm{pzc}}$, and dose. Their value here is
interpretive rather than predictive.

\textbf{Operating conditions dominate.} Permutation-importance ranks on
the hold-out set (Table~\ref{tab:ranks}) agree across six of the seven
models on the leading factors: adsorbent dose ranks first, followed by
contact time, initial concentration, and pH, with the first material
property (BET) at median rank six and the first pollutant descriptor
(log $K_{\mathrm{ow}}$) at rank nine; only the additive LR deviates.
The partial-dependence profiles of the three leading features (Fig.~S2)
show textbook adsorption behavior: removal rises steeply with dose up to
roughly 2--5 g/L and then saturates, rises within the first hours of
contact and plateaus at equilibrium, and declines with initial
concentration beyond roughly 200 mg/L as fixed capacity is exhausted.
Within known systems, the models therefore learn genuine
condition--response physics. The same ranking, however, explains the
transfer failure of Section~\ref{sec:rq2} from a second direction:
material and pollutant descriptors, the only information available about
an unseen system, contribute comparatively little to the fitted models.

\textbf{Mechanistic reading.} The SHAP dependence of
pH$-$pH$_{\mathrm{pzc}}$ (Fig.~\ref{fig:shap}a) is not smoothly
monotone but threshold-like: contributions stay within a few points
while the surface remains positively charged, and a sharp penalty of up
to nine points appears once solution pH exceeds pH$_{\mathrm{pzc}}$ by
more than about three units. Coloring this dependence by log
$K_{\mathrm{ow}}$ (Fig.~S3) shows the penalty concentrated among
hydrophilic compounds, consistent with electrostatic repulsion of
anionic species from a negatively charged surface while hydrophobic
partitioning buffers the loss for lipophilic ones. Log
$K_{\mathrm{ow}}$ itself is the strongest chemical attribution
(Fig.~\ref{fig:shap}b), monotone from roughly $-5$ points for the most
hydrophilic compounds to $+15$ for the most lipophilic. BET surface
area (Fig.~\ref{fig:shap}c) is a cautionary counterexample: the model
assigns positive contributions to records with low measured BET even
though the raw association runs the other way (median RR\% of 44 below
20 m$^2$/g versus 50 above), so this conditional attribution reflects
the correlation structure of the compilation rather than a
surface-area mechanism, and we do not interpret it mechanistically.
Attributions for the composition-averaged elemental descriptors are
small, consistent with the finding that current composition descriptors
carry little transferable information between studies. [MECHANISM
PARAGRAPH FLAGGED FOR DOMAIN REVIEW BY THE AUTHORS.]

\begin{table}[t]
\centering
\caption{Permutation-importance ranks of the leading descriptors across
the seven models (hold-out test set)
[VALUES FROM perm\_ranks\_all\_models.csv]}
\label{tab:ranks}
\begin{tabular}{lccccccc}
\toprule
Descriptor & LR & KNN & ANN & Cubist & RF & XGB & HistGB \\
\midrule
AD (dose)            & 27 & 1 & 29 & 1 & 1 & 1 & 1 \\
Time                 & 19 & 3 & 1  & 2 & 2 & 2 & 2 \\
$C_i$                & 22 & 4 & 14 & 3 & 3 & 4 & 3 \\
pH                   & 25 & 2 & 4  & 5 & 7 & 5 & 4 \\
BET                  & 24 & 5 & 2  & 7 & 6 & 6 & 6 \\
pH$_{\mathrm{pzc}}$  & 26 & 7 & 19 & 9 & 4 & 7 & 5 \\
$d_{003}$            & 8  & 6 & 8  & 8 & 12 & 8 & 10 \\
Avg.\ electronegativity & 15 & 13 & 7 & 15 & 9 & 9 & 8 \\
\bottomrule
\end{tabular}
\end{table}

\begin{figure}[t]
\centering
\includegraphics[width=\linewidth]{fig_shap_mechanism}
\caption{SHAP dependence of mechanism-linked features}
\label{fig:shap}
\end{figure}

\subsection{Computation time}
\label{sec:timing}

Training cost is not a differentiator at this data scale. On the
reference split, a single fit takes below one second for every model
except XGBoost (7--10~s across repeated measurements with the tuned
495-tree configuration), and prediction takes below 0.1~s throughout
(Table~S3). The full
leave-one-study-out audit, 59 refits of all seven models, completes
within minutes on a workstation, so the honest grouped protocol carries
no meaningful computational premium over the random-split protocol it
replaces. Hyperparameter search with Optuna dominates the overall
budget and is performed once, not per evaluation task.

\subsection{Environmental implications and outlook}
\label{sec:implications}

For practice, the results delimit two use cases sharply. Water-treatment
engineers optimizing operating conditions for an LDH--pollutant system
that has been measured before can rely on literature-trained models with
a mean absolute error near 6 percentage points of RR\%, tightened
further by conformal intervals. Materials researchers seeking to screen
new LDH formulations or new pollutants face a harder trade: the models
do carry transferable signal, reducing naive baseline errors by roughly
40 percent, but per-study skill remains near zero, so predictions for
unmeasured systems support coarse prioritization at best and should be
treated as hypotheses for experiment rather than estimates. Reported $R^2$ values
above 0.9 in prior compilations \citep{TODO-ml-adsorption-1,
TODO-biochar-cubist} are consistent with our record-level numbers and,
by the mechanism quantified here, likely overstate screening ability in
the same way; re-evaluation of published models under grouped protocols
would clarify how general this pattern is.

We see three concerns that temper the conclusions. First, the
compilation covers 59 studies, so the cross-study estimates themselves
carry sampling uncertainty; the per-study distribution, not a single
mean, is the honest summary. The second concern is reporting
heterogeneity between laboratories (differing protocols for
pH$_{\mathrm{pzc}}$ or BET), which acts as irreducible noise under (T2)
and may partly explain the transfer gap. The third concern is that our
descriptors, like those of prior compilations, describe the bulk
material; interlayer anion identity and surface morphology are rarely
reported and may carry the transferable signal that current models lack.
In future work, richer composition-structure descriptors, study-level
random effects, and few-shot calibration on a handful of measurements of
a new system are the most promising routes toward genuine screening
capability.
